\documentclass[11pt,a4paper]{article}
\usepackage[utf8]{inputenc}
\usepackage[T1]{fontenc}
\usepackage{geometry}
\geometry{left=2.7cm,right=2.7cm,top=2.6cm,bottom=2.6cm}
\usepackage{lmodern,microtype,setspace}
\usepackage{graphicx,float,amsmath,amssymb,cite,listings,xcolor,booktabs,array,caption,subcaption,algorithm,algpseudocode,tikz,pgfplots,seqsplit}
\usepackage[hidelinks]{hyperref}
\pgfplotsset{compat=1.18}
\usetikzlibrary{shapes,arrows,positioning,fit,backgrounds,matrix,calc}
\onehalfspacing
\setlength{\parindent}{1.5em}
\setlength{\parskip}{0.15em}
\hypersetup{
  pdftitle={Privacy Lens: An Evidence-Bounded Android Privacy Review Framework, Revision 24},
  pdfauthor={Zhiheng Zhang}
}
\title{Privacy Lens: An Evidence-Bounded Android Privacy Review Framework with a Signed Monotonic Legal-Review Trust Store (Revision 24)}
\author{Zhiheng Zhang \\ Student ID: 54560484 \\ Supervisor: Richard Sinnott}
\date{August 2026}
\begin{document}
\maketitle
\begin{abstract}
\small\singlespacing
Mobile privacy-review tools face an evidence-to-decision gap: a permission observation can indicate activity but cannot by itself establish purpose, necessity, lawful basis, acquisition completeness, or legal compliance. This thesis addresses that gap through Privacy Lens, a local-first Android prototype that converts bounded permission-audit summaries into traceable prompts for human review.

Following a design-science method, the architecture treats provenance, temporal scope, missing context, and uncertainty as first-class evidence rather than converting a count into a legal verdict. A fail-closed engine rejects malformed records and unsafe temporal aggregation, isolates simulator data, and records the rule-pack version applied to each finding. Jurisdiction mappings remain outside the Android, repository, and interface layers through a replaceable pack contract; the European Union General Data Protection Regulation provides the evaluated legal objective.

Revision 24 adds a canonical Ed25519-signed trust-store envelope for legal-review keys. A unique offline root signs an independently approved payload containing reviewer anchors, explicit revocations, validity dates, a safe monotonic sequence, and the accepted predecessor digest. The application persists the highest accepted sequence and digest and rejects signature mutation, unknown or malformed roots, freeze through expiry, rollback, same-sequence forks, gaps, missing history, and predecessor mismatch. The legal pack deliberately remains unattested, its verified source bytes are not bundled, and the production root and envelope stores remain empty. Findings remain local; backup and remote updates are disabled; the evaluated package has no network or sensitive runtime permission; and users can clear ordinary evidence without clearing rollback state.

In a fixed-seed 1,000-case corpus, the engine matched its specified oracle with 800 true positives and 200 true negatives. A further 1,800 independently calculated pack cases, source-content checks, signed trust-store adversarial probes, accessibility contracts, release-privacy contracts, compilation, and lint passed. The 1.12.0 APK and AAB built successfully. On one OPPO PERM00, the interface exposed official document bytes, independent review signature, and production key governance as separate links; the production trust store remained visibly unprovisioned; three short relaunches returned a live process; and the crash buffer remained empty.

These results support a stronger store-submission engineering candidate within the recorded scope. They do not establish reviewer identity or qualification, legal correctness, unrestricted Android visibility, accessibility conformance, participant outcomes, long-run reliability, production root custody, transparency, recovery across uninstall or data clearing, store approval, or publication acceptance. NIST key-lifecycle and TUF rollback concepts inform the design, but Privacy Lens claims neither a GDPR obligation derived from those sources nor TUF conformance. The contribution is an evidence-bounded architecture that separates statutory motivation, engineering controls, observed evidence, missing legal facts, and prohibited conclusions.
\end{abstract}
\newpage
\tableofcontents
\newpage
\input{Iteration-V2/Chapter1/chapter-01-23}

\subsection{Revision 24 contribution: a signed monotonic trust-store boundary}

Revision 24 replaces the preceding static reviewer-key denylist concept with a canonical, Ed25519-signed trust-store envelope. The envelope identifies an offline signing root, separates issuance from approval, carries reviewer anchors and explicit revocations, expires, advances through a positive safe-integer sequence, and binds every non-initial version to the digest of its accepted predecessor. The application persists the highest accepted sequence and envelope digest and refuses older, forked, skipped, or unchained candidates. Production roots and envelopes remain empty, so the evaluated EU pack still cannot produce legal reassurance.

The contribution is deliberately narrower than production public-key infrastructure. It demonstrates a fail-closed, testable boundary between application code and reviewer-key governance. It does not establish qualified legal review, secure root custody, witnessed publication, hardware-backed monotonicity, or continuity after uninstall or application-data clearing. This distinction continues the thesis's central method: technical evidence must make unsupported authority difficult to imply.

\input{Iteration-V2/Chapter2/chapter-02-23}

\subsection{Revision 24: key lifecycle and rollback concepts}

Key-management guidance treats cryptographic keys as lifecycle-governed objects rather than timeless constants. NIST SP 800-57 Part 1 Revision 5 discusses general key-management practices and protection across lifecycle phases \cite{nist80057r5}. Privacy Lens uses that source only as security-engineering context: neither NIST guidance nor the project's separation of issuer and approver creates a GDPR requirement or certifies a reviewer.

The Update Framework specifies version, expiry, trusted-root, and rollback-resistance concepts for software-update metadata \cite{tufspec}. Revision 24 borrows the narrower ideas of signed metadata, expiry, monotonic versions, and predecessor continuity, but does not implement TUF roles, threshold signatures, repository metadata, mirrors, consistent snapshots, or a transparency service. Calling the envelope ``TUF compliant'' would therefore exceed the evidence. The academically relevant point is the explicit translation and limitation of external security concepts, not terminological association.

\input{Iteration-V2/Chapter3/chapter-03-23}

\subsection{Revision 24 requirements and trust-store threat model}

Revision 24 adds eight requirements. \textbf{TS1} requires a fixed canonical schema whose anchor and revocation ordering cannot alter the signed payload. \textbf{TS2} requires one structurally valid Ed25519 offline root that is valid both at envelope issuance and assessment. \textbf{TS3} forbids an offline-root identifier from also identifying a reviewer anchor and requires different issuer and approver identities. \textbf{TS4} requires explicit, dated revocations within the envelope validity window. \textbf{TS5} requires sequence one to have no predecessor and later sequences to advance exactly one step while binding the accepted predecessor digest. \textbf{TS6} rejects rollback, same-sequence forks, gaps, missing history, invalid signatures, freeze through expiry, and future activation. \textbf{TS7} persists the highest accepted sequence and digest before exposing reviewer anchors and removes all effective anchors if persistence fails. \textbf{TS8} requires the interface to expose both the production provisioning state and local-history limitations.

The attacker model includes mutation of signed fields, substitution of an unknown or malformed root, replay of an older valid envelope, same-sequence replacement, omission of intermediate envelopes, predecessor substitution, and continued use after expiry. The design does not resist an attacker who can replace application code or roots, compromise the device, erase application storage, control the release channel, or subvert the offline ceremony. Those are residual distribution and platform threats, not cases silently counted as passing.

\input{Iteration-V2/Chapter4/chapter-04-23}

\subsection{Revision 24 implementation: canonical envelopes and monotonic state}

The implementation serialises a fixed field set, sorts anchors and revocations by key identifier, normalises predecessor digests, and signs the resulting UTF-8 payload with Ed25519. Structural validation checks real calendar dates, safe sequences, exact decoded key and signature lengths, identity separation, duplicate identifiers, validity intervals, and coherent revocation dates. Assessment accepts exactly one matching offline root, checks that it was valid at issuance and remains valid at assessment, verifies the signature, then evaluates stored history before it can expose reviewer anchors.

A repository stores only the highest accepted sequence and SHA-256 envelope digest under a dedicated AsyncStorage key. A pure transition validator rejects decreasing sequences, same-sequence replacement, and skipped steps; the repository fails closed to an invalid assessment with zero anchors when stored state is malformed or cannot be persisted. Clearing ordinary findings does not clear this state. The same boundary is shown in Settings and attached to finding evidence. Because AsyncStorage is app-private rather than tamper-resistant hardware, uninstall, application-data clearing, unsuitable restoration, or device compromise can break continuity. Disabled Android backup narrows one path but does not create a public transparency log.

\input{Iteration-V2/Chapter5/chapter-05-23}

\subsection{Revision 24 evaluation: signed trust-store adversarial probes}

The unchanged seeded compliance corpus returned 800 true positives, 200 true negatives, zero false positives, and zero false negatives. The 1,800-case independent governance oracle also passed. Added probes covered canonical order invariance; signature mutation; unknown, duplicate, malformed, revoked, and not-yet-valid roots; root validity at issuance; envelope expiry and future activation; issuer self-approval; root/reviewer identifier collision; incoherent post-expiry revocation; exact sequence-two advancement; replay of sequence one; a same-sequence fork; predecessor mismatch; sequence gap; missing history; idempotent reassessment; and independent rollback-state transition rejection. Application compilation, accessibility contracts across five files and 12 touchables, release-privacy contracts across 30 source files, and lint passed.

The release build produced a 62,929,891-byte APK with SHA-256 \texttt{\seqsplit{0b723dc0541b93bc27a4ffa5513d1b21f0794d39eaa49a2b719dddd5c7491097}} and a 31,790,208-byte AAB with SHA-256 \texttt{\seqsplit{ccf2e229344941d484db5e582fe8be44940f375641258ae9b7d63d4e74d3f8cf}}. Installation on one OPPO PERM00 reported version code 13, version name 1.12.0, minimum SDK 24, and target SDK 36. Three short force-stop and relaunch rounds returned a live process and the crash buffer remained empty. Device screenshots showed the legal-review warning, the three-link evidence chain, the unprovisioned production trust store, and the local-history boundary without observed clipping in the captured viewport.

These results establish conformance to the stated engineering contract only. The positive cryptographic path uses deterministic synthetic keys; no production root, signed envelope, qualified review, participant outcome, or external replication is claimed.

\input{Iteration-V2/Chapter6/chapter-06-23}

\subsection{Revision 24 residual validity and governance limits}

Rollback resistance is conditional on retained client state. An uninstall, application-data clearing event, unsuitable restoration, or compromised device can remove or alter the highest accepted sequence and digest. The application also distributes its currently empty roots and envelope through the release itself; it has no independent transparency witness, online status, threshold root control, or recovery protocol. Consequently, the mechanism raises the cost of accidental or local replay inside the admitted execution boundary but does not establish globally consistent trust metadata.

Cryptographic validity remains distinct from legal validity. A valid envelope would establish only that configured key material authorised a particular reviewer-key set. It would not establish the human signer's identity, professional qualification, independence, completeness of review, correctness of source interpretation, or continued applicability to a deployment context. Likewise, the single-device screenshots and three relaunches show a short implemented path, not general usability, comprehension, accessibility, reliability, or OEM coverage. These limitations reserve the main remaining publication-quality claims for external legal, participant, assistive-technology, multi-device, and replication evidence.

\input{Iteration-V2/Chapter7/chapter-07-23}

\subsection{Revision 24 product-position update}

Under the strict publication-target rubric introduced after the 70/100 predecessor baseline, Revision 23 was assessed at 94/100 and Revision 24 at 95/100. The one-point increase reflects stronger software and security evidence for reviewer-key governance, not a new legal-validity claim. App 1.12.0 uses version code 13, built as APK and AAB, installed on the recorded OPPO handset, and exposed the unprovisioned production trust store and rollback-state boundary. The score remains an internal expert proxy rather than a university grade, journal decision, legal opinion, accessibility certification, or participant outcome. Production ceremony, qualified review, transparent or witnessed distribution evidence, recovery, and external human studies remain gates rather than awarded credit.

\paragraph{Revision-24 conclusion.} Privacy Lens now validates a canonical Ed25519-signed reviewer trust-store envelope, requires a unique offline root valid at issuance and assessment, represents explicit revocations, and persists a monotonic sequence-and-predecessor chain. Mutation, freeze, rollback, fork, gap, missing-history, malformed-root, and persistence-transition probes fail closed. The evaluated production root list and envelope remain empty, so this additional mechanism cannot convert the unattested EU pack into reassurance. Release and bounded device evidence support the implementation claim, while local-state loss, release-channel control, qualified legal review, human comprehension, accessibility, and external replication remain unresolved. The next legal-priority step is an independently controlled production ceremony with witnessed receipts or transparency evidence and recovery semantics, followed by qualified review and participant evaluation.

{\small
\bibliographystyle{ieeetr}
\bibliography{reference}
}
\end{document}
